\section{Verification}
\label{sec:verification}

% Owner: @Capheuss (parent). Only the owner edits this file; send input to them on #40. Shared framing only; per-track text is in 5a-5d.
% Each paragraph maps to the GitHub issues noted in its comment; the
% milestone gates (M2-M4) are in issue #17 and project-docs/PLAN.md.

Section~\ref{sec:approach} describes how we measure each configuration; this
section describes how we establish that a configuration is worth measuring,
that is, that it still computes the right results and that it actually closes
the leak. STT only ever \emph{delays} instructions; it never changes what they
compute. Every check below therefore confirms that a protected core produces
exactly the architectural results of the baseline, and differs from it only in
timing. We verify in layers, from single modules up to the full SoC, so that a
bug is caught at the smallest level where it can appear. Each check is
reported only for the configurations it actually ran on; a module test does
not stand in for a full-core run, and a test that times out counts as a
failure.

% Owner: @Capheuss (track A). Only the owner edits this file; send input to them on #40. Module tests for taint tracking and the Age Matrix.

% Issues #4 (A2), #5 (A3).
\paragraph{Module tests}
The taint tables, the VP and the Age Matrix were each tested in isolation
using Chisel's built-in simulator before being wired into the core. Directed
tests of the taint tables cover a load whose destination takes its own ROB
index while its input stays untainted, the same-cycle bypass, a stalled lane
that must neither write nor bypass, the youngest source under ROB wraparound,
\texttt{x0}, third operands and moves between the integer and floating-point
registers, the strict VP boundary, taint surviving a mispredict, and a full
flush. A sweep over six register counts, including the Small to Mega sizes, and
counts that are not powers of two, checks that every register maps to its
own entry. VP tests cover out-of-order branch resolution, reuse of a branch
tag in the same cycle, wraparound and killed entries. All 30 tests pass along with
checks that the tests fail on real faults such as disabling the bypass,
truncating register indices, removing the per-cycle clearing, or removing the
Age Matrix's same-group root rule.

\paragraph{Age Matrix equivalence}
The Age Matrix was checked by equivalence using a random test which drives the 
bypass and the Age Matrix with identical inputs every cycle and compares every 
dispatched lane's YRoT, and every stored table entry, with a sequential reference model.
Its stimulus follows the method of $\mu$STT's test generator~\cite{Chen:uSTT:2025} 
and consists of dense register reuse, renaming against a live map table and free list, 
and state carried from cycle to cycle. It also adds stalls, flushes, out-of-order 
branch resolution and ROB wraparound. 
At Medium's width of 2, three seeds of 3,000 cycles compare 9,953
lanes, 441 of which read a lane dispatched in the same cycle; at $\mu$STT's
default width of 8, they compare 36,825 lanes, 3,235 of them same-cycle. We
then replayed vectors produced by $\mu$STT's own unmodified generator, 3,000
lanes in 10-wide groups, through both of our implementations. Its vectors tag
a load with a branch epoch instead of the load's ROB index, so each of our
tags is mapped to the epoch $\mu$STT assigns that load. Both implementations
match $\mu$STT's expected YRoT on every lane.

\paragraph{Full core}
With STT on, the Small and Medium cores pass all 107 RISC-V ISA tests and all
14 riscv-tests benchmarks in Verilator, as does Medium with the Age Matrix,
with no timeouts. However, at width 1, Small has no same-cycle group for the Age Matrix
to compute so it was not tested. With every switch off, the generated Verilog for Small and Medium
(591 and 594 files) matches the baseline's.
% Owner: @SamirJoshiGit (track B). Only the owner edits this file; send input to them on #40. Tests for transmitter delay and the LSU assertions.

% Issue #7 (B2): assertions at every transmitter fire point.
\paragraph{Transmitter assertions}
Because the attacks cover only a few gadgets, 
we also add a simulation assertion at every point where a transmitter can fire, 
including the load/store unit's retry, replay and forwarding paths. If a transmitter 
fires while its YRoT is not yet safe, simulation stops. This checks the STT rule on every 
instruction of every test, not only on the attack code.

% Owner: @vinle0 (track C). Only the owner edits this file; send input to them on #40. Functional regression and the attack tests.

% Issue #10 (C2): riscv-tests and Spike co-simulation in the regression script.
\paragraph{Functional regression}
Every configuration in Table~\ref{tab:configs} must pass the RISC-V ISA tests
and the riscv-tests benchmarks in Verilator, as the baseline does. Because a
delay mechanism can fail by stalling forever rather than by computing a wrong
value, a test that times out counts as a failure. This is crucial for catching
deadlocks such as a transmitter that waits on a taint that is never cleared. 
We also run Chipyard's co-simulation mode, which checks each committed instruction
against the Spike ISA simulator in lockstep, so a divergence is reported at
the first wrong instruction rather than at the end of a test. With every new
switch turned off, the generated Verilog must be identical to the baseline's,
which we check by diffing it; this confirms that our changes leave the
unmodified design untouched and baseline numbers stay comparable.

% Issues #9 (C1), #6 (B1).
\paragraph{Security checks}
The Spectre proofs of concept are the end-to-end security tests which all print
whether they recovered the secret and enable the regression script to check the
result automatically. An attack must leak on the baseline, which shows the
test itself works, and must fail on every protected configuration. Since the
conservative stand-in is strictly stronger than STT, it gives a second reference point 
such that any attack that succeeds against it points to a path that bypasses the issue queue, 
rather than to a flaw in taint tracking.

% Owner: @erickan2002-del (track D). Only the owner edits this file; send input to them on #40. Checks on the synthesized design.

% Issues #14 (D3), #15 (D4).
\paragraph{Synthesized design}
Finally, short targeted programs, such as the Spectre gadget, are run in Vivado
simulation of the synthesized design to confirm that the logic Vivado
produces behaves like the Chisel we verified in Verilator. Before that, the
testbenches released with $\mu$STT's standalone RTL are run in the same
simulator, which should reproduce the published results and validate our
simulation setup.


\paragraph{Limitations}
These checks are tests, not proofs. Passing them shows that the known Spectre
variants are blocked and that the assertions never fire on our workloads, but
not that no leak exists. Formal analysis of the RTL, as done in
SecureBOOM~\cite{Jauch:SecureBOOM:2023}, would give that stronger guarantee,
but is outside the scope of this project. The module tests check the
tag-lifetime rule (every stored taint names an instruction still in the ROB)
on the taint tables; copies of a YRoT held in the issue queues and the
load/store unit are covered by the end-to-end tests and the transmitter
assertions, not by a check of their own.
